\section{Motivação e Justificativa}

O problema central abordado por este trabalho é a detecção de faltas de alta impedância (FAIs) originadas pelo contato entre condutores de redes de distribuição e vegetação. Esse tipo de falta produz correntes de baixa magnitude insuficientes para atuar as proteções convencionais por sobrecorrente, de modo que o condutor permanece energizado indefinidamente \cite{1,2}. Em contato com material vegetal, mesmo correntes de brevíssima duração são capazes de carbonizar galhos e gerar brasas incandescentes, configurando risco direto de ignição \cite{8,Ozansoy2020}.

A importância do problema fica evidenciada pelos incêndios de *Black Saturday*, ocorridos em fevereiro de 2009 no estado de Victoria, Austrália, que resultaram em 173 mortes e na destruição de mais de 2.000 residências. Investigações posteriores apontaram falhas em ativos elétricos entre as causas de alguns focos \cite{vic_bushfire_report}. Esse episódio exemplifica um fenômeno recorrente em países com clima seco e redes aéreas extensas: a combinação de vegetação densa, condições climáticas adversas e FAIs não detectadas eleva substancialmente o risco de incêndios de grandes proporções \cite{9,vic_bushfire_report}.

Apesar das décadas de pesquisa dedicadas ao problema, ainda não existe uma solução padronizada e universalmente aceita para a detecção de FAIs \cite{Sedighizadeh2010,Aljohani2020}. Um dos principais gargalos apontados na literatura é a escassez de bases de dados reais que permitam a validação objetiva de algoritmos \cite{Gomes2021VeHIF}. A disponibilização pública do conjunto VeHIF \cite{Gomes2021VeHIF}, com mais de 900 registros de FAIs em vegetação obtidos em rede real, cria uma oportunidade concreta para avanços nessa direção.

Os benefícios esperados com a realização deste trabalho são três. Primeiro, a caracterização sistemática dos atributos dos sinais VeHIF contribui para o entendimento das assinaturas elétricas das FAIs em vegetação, ainda pouco exploradas de forma comparativa na literatura. Segundo, o conjunto de atributos extraído e avaliado pode servir de referência para o desenvolvimento de algoritmos de detecção em tempo real. Terceiro, os resultados obtidos com dados reais aumentam a validade empírica das conclusões em relação aos trabalhos baseados exclusivamente em simulação.

\section{Objetivos}

\subsection{Objetivo Geral}

O objetivo geral do trabalho é caracterizar e extrair atributos de sinais de faltas de alta impedância em vegetação a partir do conjunto de dados, avaliando sua capacidade discriminativa em um problema de classificação supervisionada.

\subsection{Objetivos Específicos}

\begin{enumerate}
\item Analisar os sinais de corrente e tensão dos canais LF e HF do conjunto VeHIF, identificando padrões visuais e estatísticos associados aos diferentes tipos de evento;

\item Definir as classes do problema de classificação com base na estrutura do dataset e nos critérios de risco de ignição estabelecidos na literatura;

\item Extrair atributos nos domínios do tempo, da frequência e tempo-frequência

\item Selecionar os atributos mais relevantes por meio de técnicas de seleção de features e análise de importância;

\item Treinar e avaliar classificadores supervisionados com os atributos selecionados, reportando métricas de desempenho adequadas ao contexto de detecção de FAIs;

\item Discutir quais atributos se mostram mais discriminativos para o problema, relacionando os resultados com as características elétricas das FAIs em vegetação descritas na literatura.

\end{enumerate}


\section{Escopo}

O trabalho propõe a caracterização e classificação automática de faltas de alta impedância associadas à vegetação (VeHIFs) em sistemas de distribuição de energia elétrica, utilizando exclusivamente o conjunto de dados público do *Vegetation Conduction Ignition Test Program* (governo de Victoria, Austrália). Compreende a revisão bibliográfica, o pré-processamento dos sinais, a extração e avaliação de atributos nos domínios do tempo e da frequência, bem como o treinamento e a comparação de algoritmos de inteligência artificial para detecção dos eventos de falta.

Não estão no escopo deste trabalho: o desenvolvimento de hardware ou relés de proteção, a coleta de novos dados experimentais, o tratamento de FAIs em superfícies não vegetais e a localização da falta na rede.

Atualmente, não há consenso na literatura sobre quais atributos e classificadores produzem melhores resultados para VeHIFs avaliados sobre dados reais, sendo a maioria dos estudos baseada em simulações ou dados laboratoriais de escopo limitado. Ao final deste trabalho, espera-se dispor de um pipeline de classificação validado e reproduzível, acompanhado de uma análise quantitativa dos atributos mais discriminativos, servindo de base para futuras extensões na área de proteção inteligente de sistemas de distribuição.


\section{Revisão bibliográgica}

As faltas de alta impedância são estudadas desde os trabalhos clássicos de \citeonline{1} e \citeonline{2}, que estabeleceram a definição do problema e documentaram a incapacidade das proteções convencionais por sobrecorrente em detectá-las. Desde então, a literatura evoluiu em três grandes frentes: métodos baseados em harmônicos e inter-harmônicos \cite{3,4}, técnicas de processamento de sinais como a transformada wavelet e a transformada de Fourier de curto prazo \cite{5,6,7}, e abordagens baseadas em inteligência artificial \cite{8}.

O caso específico das faltas associadas à vegetação (VeHIFs) recebeu atenção crescente a partir dos incêndios florestais na Austrália em 2009, que motivaram a criação do Vegetation Conduction Ignition Test Program e do dataset VeHIF \cite{vic_bushfire_report}. Gomes, Ozansoy e Ulhaq \citeonline{8} utilizaram esse conjunto de dados para demonstrar que o conteúdo de alta frequência dos sinais constitui a assinatura mais discriminativa das VeHIFs, classificando os eventos com alta sensibilidade por meio de boosted decision trees. Uma revisão recente de Yang et al. \citeonline{9} consolida o estado da arte em modelagem, detecção e avaliação de risco de ignição para faltas relacionadas a árvores.

O presente trabalho se diferencia dos estudos anteriores ao adotar uma abordagem sistemática de extração e avaliação comparativa de atributos nos domínios do tempo e da frequência diretamente sobre os dados reais do dataset VeHIF, combinada com a comparação de múltiplos algoritmos de classificação, visando identificar a combinação ótima de atributos e modelos para a detecção automática de VeHIFs.


\section{Metodologia}

O trabalho é conduzido integralmente em linguagem Python, utilizando as bibliotecas \texttt{NumPy}, \texttt{SciPy} e \texttt{PyWavelets} para o pré-processamento e a extração de atributos, e \texttt{scikit-learn} e \texttt{XGBoost} para o treinamento e a avaliação dos classificadores. A escolha de Python se justifica pela ampla adoção na comunidade científica de aprendizado de máquina, pela disponibilidade de bibliotecas maduras para processamento de sinais e pela facilidade de reprodutibilidade dos experimentos. Todo o código será versionado em repositório Git.

As principais atividades são descritas a seguir.

\begin{description}

\item[Revisão bibliográfica.] Levantamento e síntese da literatura sobre FAIs, VeHIFs, técnicas de processamento de sinais e algoritmos de classificação aplicados a sistemas de distribuição.

\item[Aquisição e análise exploratória do \textit{dataset}.] Obtenção do conjunto de dados VeHIF (\textit{Vegetation Conduction Ignition Test Program}), inspeção dos registros, definição formal das classes (falta de vegetação vs.\ operação normal) e análise estatística descritiva dos sinais.

\item[Pré-processamento dos sinais.] Filtragem, normalização, segmentação por janelas e tratamento de valores ausentes ou corrompidos, garantindo consistência do conjunto de dados antes da extração de atributos.

\item[Extração e seleção de atributos.] Cálculo de métricas nos domínios do tempo (energia, assimetria, curtose, \textit{buildup}, rugosidade) e da frequência (componentes harmônicos, inter-harmônicos e conteúdo de alta frequência via transformada wavelet e STFT). Avaliação do poder discriminativo de cada atributo por meio de métricas estatísticas e técnicas de seleção de características.

\item[Treinamento e avaliação dos classificadores.] Implementação e comparação de quatro algoritmos: árvore de decisão, \textit{Support Vector Machine} (SVM), XGBoost e rede neural artificial. Os modelos serão avaliados por validação cruzada estratificada, reportando acurácia, precisão, revocação e \textit{F1-score}. A escolha desses algoritmos cobre desde modelos interpretáveis (árvore de decisão) até modelos de maior capacidade (XGBoost e redes neurais), permitindo uma comparação abrangente do \textit{trade-off} entre desempenho e complexidade.

\item[Análise dos resultados e redação final.] Discussão dos resultados obtidos, identificação dos atributos e modelos mais adequados para o problema, e elaboração do documento final do TCC.

\end{description}
